\documentclass[11pt]{article}

\usepackage[margin=1in]{geometry}
\usepackage{graphicx}
\usepackage{amsmath}
\usepackage{amssymb}
\usepackage{authblk}
\usepackage[hidelinks]{hyperref}
\usepackage{natbib}
\usepackage{booktabs}
\usepackage{caption}
\captionsetup{font=small,labelfont=bf}

\renewcommand\Authfont{\normalsize}
\renewcommand\Affilfont{\small\itshape}

\title{\bfseries A scaling law of contextual persistence in human language}

\author[1]{Elan Barenholtz}
\affil[1]{Department of Psychology and Center for Complex Systems, Florida Atlantic University}

\date{}

\begin{document}
\maketitle

\begin{abstract}
\noindent
Human language obeys scaling laws for vocabulary growth, word frequency, and syntactic
structure, yet none has been established for how the influence of prior context decays
during production. Using large language models as calibrated probes, we quantified how
strongly prior context constrains upcoming words, isolating order-specific structure with a
shuffled-token baseline; the contextual persistence function $P(d)$ measures this influence
at distance $d$. Across ten corpora (449 documents) spanning six language families and
written and spoken modalities, $P(d)$ decayed approximately as $1/d$ (mean exponent $-1.04$;
95\% CI $-1.15$ to $-0.93$; median $r^2=0.96$) across nearly two orders of magnitude, vanished
in token-scrambled and synthetic controls, and replicated across independent probes. Because
the exponent is close to $-1$, contextual influence is distributed approximately uniformly
across logarithmic timescales; its value also lies in the same order-unity range as independent
measurements of human memory decay. These results identify a scaling law of contextual persistence
that any theory of language production must explain.
\end{abstract}

\section{Introduction}

Many complex systems exhibit scale-free statistical regularities, in which no single
scale characterizes the behaviour and structure repeats across magnitudes. Human
language is a rich source of such regularities: the frequency of words follows Zipf's
law \citep{zipf1949human}, vocabulary grows with text length under Heaps' law
\citep{heaps1978information}, the length of linguistic constituents scales inversely
with the size of the whole under Menzerath's law, and fluctuations scale with the mean
under Taylor's law \citep{altmann2016statistical}. These laws are quantitative,
reproducible, and largely invariant across languages, and they have shaped how the
organization of language is understood. Statistical laws of this kind have played an
outsized role in the study of complex systems precisely because they reveal deep
constraints while remaining largely agnostic about mechanism. Zipf's law preceded any
accepted mechanistic explanation by decades, yet became a central empirical target against
which theories of language were measured; the value of such a law lies less in the
mechanism it immediately implies than in the fact that every candidate mechanism must
reproduce it. Crucially, Zipf's law is compelling not merely because word frequencies follow
a power law but because the exponent sits near unity across languages and corpora: the
reproducibility of the value, not the power-law form alone, is what marks a genuine law. Yet
one basic dimension of language has no such law: how strongly the past
constrains the future during production --- how the
influence of prior material decays with distance from the current point of generation.
Contextual dependence is central to coherent discourse, but existing characterizations
--- discourse models, situation models, resonance-based accounts --- operate
representationally rather than quantitatively, and conventional measures of statistical
dependence conflate sequential structure with vocabulary-level regularity. Whether
contextual influence in language obeys a scaling law, and if so of what form, is
unknown.

Here we measure this decay directly. The contextual influence carried by the past is
the degree to which knowing the preceding words improves prediction of what comes next,
beyond what the words' identities alone would afford. We operationalize this with large
language models as calibrated probes of conditional probability: for a fixed upcoming
span, we ask how much revealing the true, ordered prior context reduces the probe's
surprise, and subtract the reduction obtained when the same words are supplied in
scrambled order. The difference removes any benefit attributable to vocabulary or topic
alone and retains only what depends on sequential arrangement --- the order-specific
contextual influence. Measured as a function of distance $d$ back from the point of
production, this defines the \emph{contextual persistence function} (CPF), $P(d)$; its
functional form we call the \emph{contextual persistence law}. We apply it across ten
corpora spanning six language families, multiple genres, and written and spoken modalities.

$P(d)$ decays approximately as $1/d$ with mean slope $-1.04$ (SD $=0.15$) and no
characteristic cutoff across nearly two orders of magnitude in distance. This decay
is consistent across language families, genres, and modalities, vanishes in
random-token controls, scales with the integrity of sequential structure, and is
broadly distributed across prior sentences rather than concentrated in a small number
of anchors. The CPF describes a continuous, scale-free temporal
structure in human language, measurable directly from the linguistic sequence and
replicable across independently trained probe models
(Llama-3.1-8B, Mistral-7B).

The shape of $P(d)$ places language alongside other human behaviours and neural
processes that exhibit scale-free temporal structure, and constrains the class of
mechanisms capable of producing such structure. In particular, the absence of any
characteristic scale across the measured range --- combined with the graded
contributions of prior content distributed across the sequence --- is difficult to
reconcile with architectural accounts of memory that posit a sharp division between
maintained and retrieved material with distinguishable temporal dynamics. The pattern
is consistent instead with mechanisms in which prior content continues to shape
ongoing generation through the evolving probability structure of the sequence itself.

\section{Results}

\subsection{The contextual persistence function}

We define the contextual persistence function (CPF), $P(d)$, as the predictive influence of prior
context at distance $d$ from the current point of generation, estimated as the
reduction in target perplexity produced by revealing prior context at that distance,
with shuffled-token baselines controlling for the contribution of token identity
alone (Methods). $P(d)$ provides a distance-resolved measurement of how preceding
language constrains the probability structure of upcoming tokens.

Across all corpora, $P(d)$ is positive at distances well beyond any plausible bound
on working-memory capacity. At distances approaching 1000 tokens (the far end of our
measurement range), prior context continues to reduce target perplexity by a measurable
margin relative to shuffled controls. Long-range contextual influence is therefore not confined to a recency
window but extends as a graded function across the full sequences we examined.

\subsection{Heavy-tailed decay across languages, genres, and modalities}

The functional form of $P(d)$ is well approximated by a heavy-tailed decay. Figure~1
plots $P(d)$ against distance on log--log axes for each of the ten corpora. Across
languages, genres, and modalities, $P(d)$ decays smoothly across nearly two orders of
magnitude with a mean slope of $-1.04$ (SD $=0.15$) --- close to a $1/d$ power law
--- and a median goodness-of-fit of $r^2 = 0.96$. Fits are computed on $d \geq 10$
tokens, the range over which both ordered- and shuffled-context bases are stable;
very-short-range intervals are reported in Methods. In a comparison of five candidate
functional forms by small-sample-corrected information criteria (Supplementary
Information), the power law is preferred in eight of the ten corpora and outperforms a
pure exponential in all ten --- a fixed-scale exponential is never preferred. The two
remaining corpora, Japanese literary prose and Buckeye spoken English, show mild downward
curvature and are marginally better described by a log-normal and a stretched-exponential
form respectively; in both the decay remains heavy-tailed and extends well beyond the range
over which any bounded buffer could plausibly contribute.

The cross-corpus consistency of the decay regime is notable. The sample spans
Germanic, Romance, Turkic, Japonic, Uralic, and Slavic language families, as well as
fiction and literary prose, news and expository prose, prepared speech transcripts,
and spontaneous speech. A narrow heavy-tailed regime across this range indicates that
the CPF reflects a general property of natural language rather than a
property of any single language, register, or modality.

The qualitative pattern that defines $P(d)$ --- a positive contribution of ordered
context at long range together with a near-zero contribution of shuffled context ---
extends beyond the six families covered by the long-range run. In a complementary
cross-language analysis evaluating shorter context lengths on both probe models, the
long-range sign pattern is preserved across three additional language families:
Afro-Asiatic, Koreanic, and Sino-Tibetan (Supplementary Information). The headline
slope of $P(d)$ is therefore obtained on a six-family set within a broader nine-family
cross-language replication of the qualitative pattern that defines the persistence
function.

\subsection{An exponent near $-1$: uniform influence across logarithmic timescales}

Two properties of the decay should be separated. The first is that it is scale-free: a power
law has no distance at which its behaviour changes character, and our fits show none across
the measured range. The second, and more distinctive, is the value of the exponent. Across the
ten corpora the fitted exponent is close to $-1$ and clustered around it (mean $-1.04$, SD $=0.15$).
Document-level bootstrap (Supplementary Information) shows this between-corpus spread modestly
exceeds within-corpus sampling error (typical bootstrap SE ${\approx}\,0.10$), indicating small but
genuine cross-corpus heterogeneity around a central value near $-1$ rather than a single exact
constant. The reproducibility of the near-unity value across otherwise very different samples is
itself part of the result, and it carries a specific quantitative meaning.

For a CPF of the form $P(d) \propto d^{-\alpha}$, the aggregate contextual influence contained
within a logarithmic interval of distance --- the band spanning one multiplicative step, or one
decade, of the past --- is proportional to $d\,P(d) \propto d^{\,1-\alpha}$. When $\alpha \approx 1$
this quantity is approximately constant. Put concretely, if one divides the past into logarithmic
bands --- 1--10, 10--100, and 100--1{,}000 tokens back --- each band contributes approximately the
same aggregate contextual influence, even though the influence of any single word continues to
decline with distance. Contextual influence is therefore distributed approximately uniformly across
logarithmic timescales: no logarithmic timescale dominates the aggregate, and every decade of the
past contributes comparably to ongoing prediction.

This should not be read as memory being equally strong at all distances --- individual words do
become progressively less influential as they recede. What is approximately equal is the
\emph{aggregate} influence summed within each logarithmic decade. The distinction is precise:
pointwise influence falls as $1/d$, while cumulative influence per multiplicative interval of
distance remains flat.

The value $\alpha \approx 1$ is also a boundary. For $\alpha > 1$, aggregate influence is
concentrated near the present and the remote past contributes negligibly in sum; for $\alpha < 1$,
the sum is increasingly dominated by distant context. $\alpha \approx 1$ is the crossover between
these regimes --- the exponent at which no timescale dominates the total. That human language sits
at this boundary, rather than at an arbitrary point along the continuum, is what most distinguishes
the result from a generic heavy tail.

The near-constancy of the exponent is what makes this a law rather than a curve fit. As with
Zipf's law, the reproducibility of the value --- a quantitative constant near $-1$ recurring
across languages, modalities, and independently trained probe models --- is the phenomenon, not
merely the power-law form.

\subsection{The CPF reflects sequential structure, not properties of the probe}

A central concern in any LLM-based measurement is whether the observed pattern
reflects properties of natural language or properties of the probe model. We address
this directly with synthetic-sequence controls. We constructed sequences with
token-frequency statistics matched to natural text but with sequential structure
removed. Raw perplexity over these synthetic sequences exhibits apparent power-law
structure with a slope comparable in magnitude to that observed on natural text
(Llama-3.1: $\approx -1.27$ on uniform random vocabulary; Mistral-7B: $\approx -1.34$),
replicating the general tendency for language-model perplexity to fall as context
length increases even when genuine long-range dependence is absent.

The order-specific CPF, $P(d)$, however, is near zero across all
distances on the synthetic sequences. The shuffled-token baseline that defines $P(d)$
absorbs the perplexity reduction attributable to token statistics alone, isolating the
contribution of sequential organization. The long-range, heavy-tailed $P(d)$ observed
for natural language therefore depends on properties of the linguistic sequence
itself, not on the probe. We replicated the persistence signature with two
independently trained probe models, Llama-3.1-8B and Mistral-7B, obtaining the same
qualitative pattern and comparable decay ranges on shared corpora, along with near-zero
$P(d)$ on synthetic controls. The CPF is therefore not an idiosyncratic
property of the model used to estimate it.

\subsection{Order and content make distinguishable contributions to persistence}

To characterize the structure of $P(d)$ further, we compared ordered context to two
structure-disrupting manipulations. Sentence shuffling preserves within-sentence syntax
and lexical content while disrupting discourse-level ordering. Token shuffling
eliminates ordered structure within the contextual span.

Sentence shuffling reduces $P(d)$ at short distances, especially below approximately 30
tokens, but leaves much of the long-range tail intact (Figure~2). Token shuffling
collapses $P(d)$ toward zero across all distances. The CPF therefore
decomposes into two distinguishable contributions: a chaining-specific component
concentrated at short range (slope $\approx -1.26$, approaching zero by $d > 100$
tokens) and a content-driven component that extends across the long-range tail (slope
$\approx -0.78$). Discourse-level chaining shapes the short-range structure of $P(d)$;
content-driven integration carries its long-range tail.

\subsection{Persistence is broadly distributed across prior sentences}

We next asked whether persistence is concentrated in a small number of privileged prior
sentences or broadly distributed across the prior context. We addressed this with a
sentence-level ablation: for each target, we removed one prior sentence at a time and
measured the resulting change in prediction.

The majority of prior sentences contribute positively to prediction at all distances
examined. The proportion of contributing sentences declines gradually with distance,
from approximately 90\% at short range to 55--60\% beyond 500 tokens, with no sharp
transition or threshold (Figure~3). The magnitude of contribution is concentrated at
recent positions (top-5 sentences account for 30--51\% of total positive influence;
Gini coefficient 0.78--0.90), but this concentration is attributable to recency:
variation in mean influence by absolute position within the prior context is accounted
for by distance to the target, with no detectable residual effect of structural position
(e.g., document opening). Persistence is therefore a distributed field of influence across the
prior sequence, with magnitude graded smoothly by recency rather than localized to a
small number of anchoring positions.

\section{Discussion}

We have identified a quantitative regularity, which we term the
contextual persistence law: contextual influence in human language decays as an approximately
scale-free power law, with a near-constant exponent across ten corpora, six language families,
and both written and spoken modalities. Because that exponent is close to $-1$, contextual
influence is distributed approximately uniformly across logarithmic timescales rather than
concentrated at any one --- a more specific statement than the mere absence of a horizon. The
consistency of the exponent across such diverse samples makes this a candidate \emph{universal}
regularity of human language, in the sense that Zipf's and Heaps' laws are universal. We read it
as a law to be explained rather than as an adjudication between mechanisms, and its main force is
as a constraint. The evidence establishes the regularity; what remains open is why it exists, not
whether it is real. Any successful account of production must explain why the statistical product of
human language lacks a characteristic temporal scale. This is a sharper demand than it may
appear. An architecture with a fixed maintenance horizon would generally be expected to
produce an inflection in $P(d)$ near that horizon, of whatever size; we observe none across
nearly two orders of magnitude. An account in which long-range coherence rests on a small
number of privileged anchor positions would generally predict a concentrated rather than a
distributed influence field; we observe broad, recency-graded distribution. We do not take
these observations to settle how memory is organized during production --- they characterize
the product, not the process --- but the law they describe is what every candidate process
now has to reproduce.

The mean exponent we report ($-1.04$, SD $=0.15$) sits in a neighborhood that has been
independently identified in two adjacent traditions in memory research.
\citet{anderson1991reflections}, examining the recurrence statistics of information demand
in natural environments, reported a power law with exponent of approximately $-0.77$.
\citet{kahana2002age}, examining associative decay in free recall via lag-conditional
response probability, reported a forward-decay exponent of approximately $-0.82$ for
younger adults. Each of these measurements was made by entirely different methods on
entirely different data, and --- crucially --- neither involves a language model: one is a
statistical analysis of environmental information demand, the other a behavioural memory
task. That two model-free measurements of human memory land in the same order-unity exponent
range as our production-side measurement indicates that the scale-free form is neither an
artefact of our instrument nor idiosyncratic to language, but a property that recurs across
paradigms that are unambiguously about memory. The exponents are not identical --- recall and
environmental statistics sit somewhat shallower than the production exponent we report --- and
we do not claim a single mechanism generates all three. But their common scale-free form, absent
any characteristic bound in any of the three, is the pattern a graded, unbounded memory would
naturally produce, and one that a fixed-horizon architecture would not naturally predict.

This empirical picture aligns with a longer trajectory in cognitive psychology away from
architecturally bounded views of memory and toward continuous and graded characterizations.
Cowan's focus-of-attention model \citep{cowan2001magical} recasts working memory as a graded
availability of items rather than a fixed-capacity slot store; Ericsson and Kintsch's
long-term working memory framework \citep{ericsson1995longterm} argues that skilled
performers make routine use of long-term-memory structures during ongoing task performance,
blurring the maintenance-retrieval boundary; and continuous-resource accounts in visual
working memory \citep{ma2014changing} replace discrete-item architectures with graded
distributions of representational precision. The CPF provides a direct
quantitative measurement of an analogous continuity in language production, observed at the
scale of token-level sequential dependence and across naturalistic discourse rather than
experimental task sequences. To the extent the law bears on this debate, it is consistent
with that longer trajectory rather than a refutation of it.

A natural objection is that our measurement characterizes the statistical structure of
language rather than the memory process that produced it: the CPF is computed
over text, and predictive structure in text need not imply that speakers and writers actively
used long-range memory to generate it. We take this objection seriously and delimit our claim
accordingly. What the data establish is a property of the human-produced signal, not of any
particular generative mechanism. That property is, however, far more constrained than a generic
statistical regularity. Because $P(d)$ is defined by subtracting a shuffled-token baseline, it
isolates the component of contextual influence that depends on the order of prior words; the
near-zero $P(d)$ on token-scrambled and frequency-matched synthetic sequences shows it cannot
arise from vocabulary statistics or from the probe's own calibration, and its replication across
independently trained probes shows it is not an artefact of any single model. The long-range
dependence therefore reflects genuine order-dependent structure in what humans produce, not
lexical co-occurrence or measurement bias.

What the data do not establish is that production is implemented by a sequential process reading
from an unbounded memory; that mechanistic claim the present measurements cannot adjudicate, and
we do not make it. A natural question is whether autoregressive systems --- which generate each element
by predicting it from unboundedly long prior context --- reproduce the same scale-free form under
identical measurement. Establishing this would be a sufficiency result, showing that a purely
sequential generation mechanism can produce the observed structure and placing the burden on
fixed-horizon or non-sequential accounts to explain the absence of a characteristic scale by some
other route; we leave that direct comparison to future work.
Converting the present result from a property of the product into a claim about the producer
will require production-time evidence --- for instance, whether human processing costs,
disfluencies, or neural responses during generation scale with $P(d)$.

Human language appears to obey a simple quantitative law --- the contextual
persistence law --- in which the influence of prior context is distributed approximately
uniformly across logarithmic timescales during production. Whether this regularity reflects
properties of memory, discourse organization, planning, prediction, or communication is
precisely the mechanistic question the law leaves open; whether it generalizes
to non-linguistic sequential behaviours with extended temporal coherence --- music improvisation,
motor sequences, drawing --- remains to be tested. But like other scaling laws in language, its
value lies less in any mechanism it immediately reveals than in the constraint it imposes on
every mechanism proposed thereafter: its existence is now an empirical target that any theory of
language production must account for.

\section{Methods}

\subsection{Corpora}

The long-range CPF was computed on ten corpora: English fiction (Project
Gutenberg), English news prose, English spontaneous speech (Buckeye Corpus of Conversational
Speech), TED transcripts in English, German, French, Turkish, and Russian, and literary prose
in Japanese and Finnish. Together these cover six language families (Germanic, Romance, Turkic,
Slavic, Japonic, Uralic) and four production modes (written fiction, news prose, prepared
speech, spontaneous speech). Per-corpus document counts (documents contributing at least one
full-length target window) were: English fiction 60, English news 60, Buckeye 26, TED English 60,
German 60, French 31, Turkish 12, Russian 60, Japanese literary 42, Finnish literary 38 --- 449
documents in total. The complementary cross-language qualitative replication of the sign-flip
pattern (Afro-Asiatic, Koreanic, Sino-Tibetan) was performed on Wikipedia and TED-transcript cells
from a broader corpus inventory at shorter context lengths. Full corpus details, sources, and
preprocessing steps are provided in Supplementary Information.

\subsection{Probe models}

Two autoregressive language models served as probes: Llama-3.1-8B (base;
\texttt{Meta-Llama-3.1-8B}) and Mistral-7B-v0.1. Neither was fine-tuned for this study. The
probes were used only to compute conditional probability estimates over fixed token sequences;
the language models themselves were not the object of study. Headline results were obtained on
Llama-3.1-8B; the qualitative pattern was replicated on Mistral-7B for shared corpora as a
probe-independence check.

\subsection{Contextual persistence function}

For each document, a single target region of 30 tokens was selected at the midpoint (50\%) of
document length. Context preceding the target was revealed in incrementally larger spans at
log-spaced context lengths $c \in \{0, 1, 2, 4, 8, 16, 32, 64, 128, 256, 512, 1024\}$. At each
context length, target perplexity was computed under intact context (the actual preceding $c$
tokens, original order) and shuffled context (the same $c$ tokens, random order). Per-distance
marginals were obtained as the per-token change in perplexity across successive context lengths,
$m_d = (\mathrm{ppl}_{c_{i-1}} - \mathrm{ppl}_{c_i}) / (c_i - c_{i-1})$, with distance indexed by
the geometric mean $d = \sqrt{c_{i-1} c_i}$ of the interval. The CPF was defined
as $P(d) = m_d^{\text{intact}} - m_d^{\text{shuffled}}$, isolating the contribution of sequential
structure from probe-internal token-identity calibration. Marginals were aggregated across
documents within each corpus and binned by distance. Only documents with at least
$1024 + 30 + 50$ tokens were included, ensuring a full-length context window and target for every
measurement.

\subsection{Functional-form fitting}

The per-token order-specific gap was fit to five candidate functional forms --- power law,
exponential, stretched exponential, truncated power law, and quadratic-in-log
(log-normal-in-distance) --- and ranked per corpus by the small-sample-corrected Akaike
criterion (AICc); the full comparison is reported in Supplementary Information. Power-law fits
were computed on positive bins with $d \geq 10$ tokens, the first interval at which both ordered
and shuffled bases are at least eight tokens. Shorter intervals were excluded because with only one or
two context tokens there is no meaningful permutation, leaving the shuffled baseline --- and hence
the order-specific difference $P(d)$ --- unstable at $d < 10$; these intervals are neither fit nor
plotted.

\subsection{Synthetic-sequence controls}

To distinguish probe-internal calibration dynamics from genuine sequential dependence, two
synthetic-sequence controls were constructed: uniform-vocabulary sequences (tokens drawn uniformly
at random from the model vocabulary) and frequency-matched sequences (tokens drawn with empirical
token-frequency probabilities matched to the natural corpora). The same CPF
protocol was applied to both.

\subsection{Sentence-level decomposition}

The contribution of discourse-level ordering versus within-sentence content was assessed by
computing $P(d)$ under three context conditions: ordered (intact), sentence-shuffled (sentences
permuted while within-sentence token order preserved), and token-shuffled. The chaining-specific
component was estimated as the difference between ordered and sentence-shuffled persistence; the
content-driven component as the difference between sentence-shuffled and token-shuffled. This
decomposition was computed across eight cells: the long-range corpora excluding Buckeye (spontaneous
speech lacks unambiguous sentence boundaries) and Russian (added to the headline analysis after the
decomposition run). Figure~2 aggregates these eight cells.

\subsection{Sentence-level ablation}

For documents in two corpora (English fiction and English TED transcripts, 20 documents per corpus,
1024-token prior context), each prior sentence was individually removed and the resulting change in
target perplexity was computed (leave-one-sentence-out). Per-sentence influence was characterized by
sign (positive $=$ removing the sentence increased target perplexity) and magnitude. The proportion
of contributing sentences was computed in distance bins, and concentration metrics (Gini coefficient
over per-sentence magnitudes; share of total positive influence in the top-5 sentences) were computed
per document and aggregated. To test whether magnitude variation across prior-sentence positions
reflects a privileged structural role over and above distance from the target, the full
leave-one-sentence-out influence matrix was computed on a 15-document subset of each ablation corpus;
variation in influence magnitude was accounted for by distance, with no detectable residual effect of
structural position once distance was controlled.

\section*{Supplementary Information: functional-form robustness}

Because scale-free claims are sensitive to the choice of fitting procedure and to
competing functional forms, we treated functional-form selection as a robustness analysis
rather than assuming a power law. For each corpus the per-token order-specific gap was fit,
in log space, to five candidate forms: a power law ($\log y = a + b\log d$), an exponential
($\log y = a + bd$), a stretched exponential ($\log y = a - \beta d^{\gamma}$), a truncated
power law ($\log y = a + b\log d - d/\lambda$), and a quadratic-in-log (log-normal-in-distance)
form ($\log y = a + b\log d + c(\log d)^2$). Models were ranked per corpus by the small-sample
corrected Akaike criterion (AICc), which penalizes the additional free parameters of the more
flexible forms; each fit used the seven binned points with $d \geq 10$ and positive influence.

The power law is preferred in eight of the ten corpora (Table~S1) and outperforms the pure
exponential --- the canonical fixed-scale alternative --- in all ten, by $\Delta\text{AICc}$
ranging from $+0.3$ to $+16.8$ (median $\approx +14$); the exponential is never preferred. The
truncated power law never improves on the plain power law: in six corpora its fitted decay term
is non-negative (no cutoff at all), and in three of the remaining four the implied cutoff
$\lambda$ lies beyond the measured range. Only Buckeye spontaneous speech shows a truncated-power-law
cutoff within range ($\lambda \approx 276$ tokens). Separately --- and as a distinct model comparison
--- Buckeye is also the one corpus for which the stretched exponential is marginally preferred over the
power law by AICc, while Japanese literary prose marginally prefers the quadratic-in-log form. Both the
cutoff and the stretched-exponential preference point to mild downward curvature at the longest distances
in this cell, but they are separate tests and should not be read as one explaining the other. In both
non-power-law cases the power law remains a close second (margin $<2$ AICc), the fitted stretched
exponent stays well below one, and the decay remains heavy-tailed. The fitted power-law exponents
are tightly clustered near $-1$ (range $-0.75$ to $-1.31$; eight of ten within $-0.93$ to $-1.20$).
We note that the maximum-likelihood, Kolmogorov--Smirnov procedure of \citet{clauset2009power} is
designed to test whether a \emph{sample} is drawn from a power-law distribution, and is not
directly applicable to a regression of $P(d)$ on distance; the information-criterion comparison
above is the appropriate test for competing decay forms. No alternative form removes the central
observation: across the measured range the decay carries no characteristic scale.

\renewcommand{\thetable}{S1}
\begin{table}[h]
\centering
\small
\begin{tabular}{lcccc}
\toprule
Corpus & $N$ & Preferred form & $\Delta$AICc (exp $-$ power) & Exponent [95\% CI] \\
\midrule
Gutenberg (fiction, en) & 60 & power law & $+12.9$ & $-1.05$ [$-1.33$, $-0.90$] \\
TED (en)                & 60 & power law & $+15.5$ & $-0.97$ [$-1.11$, $-0.87$] \\
TED (de)                & 60 & power law & $+15.5$ & $-1.03$ [$-1.16$, $-0.93$] \\
TED (fr)                & 31 & power law & $+9.4$  & $-1.31$ [$-1.57$, $-1.02$] \\
TED (tr)                & 12 & power law & $+5.0$  & $-0.75$ [$-1.02$, $-0.37$] \\
TED (ru)                & 60 & power law & $+16.8$ & $-1.02$ [$-1.13$, $-0.94$] \\
Literary prose (ja)     & 42 & quadratic-in-log & $+13.2$ & $-0.98$ [$-1.15$, $-0.84$] \\
Literary prose (fi)     & 38 & power law & $+14.9$ & $-0.93$ [$-1.07$, $-0.77$] \\
News (en)               & 60 & power law & $+13.6$ & $-1.16$ [$-1.37$, $-0.98$] \\
Buckeye (speech, en)    & 26 & stretched exp. & $+0.3$ & $-1.20$ [$-1.53$, $-0.82$] \\
\bottomrule
\end{tabular}
\caption{Functional-form comparison and per-corpus exponents across the ten long-range corpora
(Table~S1). $N$ is the number of contributing documents. $\Delta$AICc (exp $-$ power) is positive
whenever the power law is preferred over a pure exponential; it is positive in every corpus, and the
power law is the AICc-preferred form in eight of ten (the two exceptions show mild downward curvature,
with the power law a close second). Exponent 95\% confidence intervals are from a document-level
bootstrap ($B=2000$); the wider intervals correspond to the smaller corpora (Turkish, Buckeye, French).}
\end{table}

\subsection*{Robustness of the exponent: units and heterogeneity}

Two further checks bound the interpretation of the exponent. First, because $P(d)$ is defined as a
difference of perplexities, and perplexity is exponential in the underlying log-loss, we recomputed
the entire analysis with $P(d)$ defined in log-probability (nats) units --- i.e.\ replacing each
perplexity by its logarithm, an additive information quantity --- and refit the power law. The
scale-free \emph{form} is robust to this choice: every corpus remains a near-$1/d$ power law and the
ordering of corpora is preserved. The precise exponent, however, is not unit-invariant: the mean
shifts from $-1.04$ (perplexity) to $-1.18$ (log-probability), a consistent offset of ${\approx}\,0.14$
present in every cell. We report perplexity-based values in the main text for continuity with the raw
measurement and note that the order-unity exponent, and the associated approximate uniformity of
influence across logarithmic timescales, hold in both unit conventions, with the exact value
dependent on the choice.

Second, a document-level bootstrap ($B=2000$ resamples of documents within each corpus) yields the
per-corpus confidence intervals in Table~S1. The between-corpus standard deviation of the ten point
estimates ($0.156$) modestly exceeds the typical within-corpus bootstrap standard error
(${\approx}\,0.10$); the implied genuine between-corpus component is ${\approx}\,0.12$. The exponent is
therefore best described as clustered near $-1$ with small but real cross-corpus heterogeneity rather
than as a single universal constant, with the low-$N$ corpora (Turkish, Buckeye, French) contributing
most of both the spread and the widest intervals.

\section*{Data and code availability}

Analysis code, corpus manifests, per-corpus persistence curves, and the figure-generation pipeline
will be made publicly available in a repository with an assigned DOI upon publication; they are
available from the author on reasonable request in the interim.

\section*{Funding}

The author received no specific funding for this work.

\section*{Competing interests}

The author declares no competing interests.

\begin{figure}[t]
\centering
\includegraphics[width=\textwidth]{figures/fig1_persistence.pdf}
\caption{\textbf{Cross-corpus contextual persistence functions.} Per-token order-specific contextual
influence $P(d)$ as a function of distance $d$, log--log axes. Each panel is one of ten corpora spanning
six language families (Germanic, Romance, Turkic, Slavic, Japonic, Uralic) and four production modes. All
cells exhibit heavy-tailed decay across nearly two orders of magnitude in distance, with mean power-law
slope $-1.04$ (SD $=0.15$) --- close to a $1/d$ law. Points are shown for $d \geq 10$ tokens; the dotted
line is the power-law fit. Intervals below $d = 10$ are excluded because at such short context the
shuffled-token baseline is unstable --- with only one or two tokens there is no meaningful permutation, so
the order-specific difference is ill-defined there and does not reflect the persistence decay (Methods).
The two cells marked $\dagger$ (Japanese literary prose; Buckeye spontaneous English speech) show mild
downward curvature and are marginally better fit by a log-normal and a stretched-exponential form
respectively, but remain heavy-tailed with the power law a close second.}
\end{figure}

\begin{figure}[t]
\centering
\includegraphics[width=0.85\textwidth]{figures/fig2_sentence_shuffle.pdf}
\caption{\textbf{Order and content make distinguishable contributions to persistence.} $P(d)$ under
context conditions that isolate discourse-level ordering from within-sentence content. The intact and
sentence-shuffled curves coincide at long range and diverge only below ${\sim}30$ tokens; both lie far
above the token-shuffled baseline at all distances. The CPF decomposes into a
chaining-specific component (slope $\approx -1.26$, approaching zero by $d > 100$ tokens) and a
content-driven component (slope $\approx -0.78$) that carries the long-range tail.}
\end{figure}

\begin{figure}[t]
\centering
\includegraphics[width=0.85\textwidth]{figures/fig3_distributed_influence.pdf}
\caption{\textbf{Persistence is broadly distributed across prior sentences.} Proportion of prior
sentences whose removal degrades target prediction as a function of distance to the target. The
proportion declines smoothly from ${\sim}90\%$ within the first 32 tokens to ${\sim}55$--$60\%$ beyond
500 tokens, with no sharp transition. The top-5 sentences account for 30--51\% of total positive
influence (Gini 0.78--0.90), but this concentration reflects recency, not a privileged structural role
for any specific position.}
\end{figure}

\bibliographystyle{plainnat}
\bibliography{references}

\end{document}
